% Chapter 5 of Deep Learning as an Experimental Science: controls for a search method.
\labchapter{Controls}{Random search as the placebo arm}{Neural architecture search finds better architectures than random search.}{As reported for ENAS (Pham et al.\ 2018) and DARTS (Liu, Simonyan and Yang 2019)}{Vanished: trained with the same pipeline, architectures sampled at random from the same space do about as well.}{ch:controls}
\begin{refsection}[chapters/05-controls.bib]

\begin{chaptersummary}
In 2018 and 2019 weight sharing made neural architecture search cheap, and ENAS and DARTS reported architectures much better than those of random search. When later studies gave the random arm the same space, the same training pipeline and enough draws, the searched architectures were ahead by a few tenths of a point on CIFAR-10 and not at all on Penn Treebank; the original gaps came from tiny random arms, from proxies that barely ranked architectures, and from pipelines that differed between arms. The method is the control arm of the older sciences: a placebo that samples the space, a sham that spends the same budget without the search strategy, and positive and negative controls. We prove that a search guided by a Gaussian proxy gains over the placebo the proxy's correlation with the truth times the gain of a perfect search, and reports a gain inflated by the inverse square of that correlation. On a benchmark of 192 digits networks, random search with a short-training proxy does at least as well as evolution, and a baseline copied from a weaker pipeline makes the search look more than twice as useful as it is.
\end{chaptersummary}

\section{The claim: better than random, at a thousandth of the cost}\label{sec:controls:claim}

The first architecture searches were expensive.\index{neural architecture search} \Textcite{ZophLe2017} trained a recurrent controller with reinforcement learning to write down network descriptions, with the validation accuracy of each trained network as the reward, and found a CIFAR-10 network with 3.65\% test error and a recurrent cell that beat the LSTM on Penn Treebank. Searching for a repeatable cell made the result transferable to ImageNet \parencite{ZophEtAl2018}, and evolution in the same cell space produced AmoebaNet \parencite{RealEtAl2019}, at search costs that \textcite{LiuSimonyanYang2019} list as 2000 and 3150 GPU-days.

\paragraph{Weight sharing.} \Textcite{PhamEtAl2018} cut the cost by three orders of magnitude with \term{weight sharing}\index{weight sharing}. In ENAS every candidate architecture is a subgraph of one large graph and inherits its weights, so no candidate is trained from scratch: a controller samples subgraphs and is trained by policy gradient to maximize a reward measured on validation data, while the shared weights are trained on the sampled subgraphs. On one GTX 1080Ti the search took less than sixteen hours, and ENAS found a recurrent cell with test perplexity 56.3 on Penn Treebank and a convolutional cell with 2.89\% test error on CIFAR-10 when trained with cutout, 3.54\% without.\sidenote{The arXiv version of the paper reports a perplexity of 55.8; the proceedings of ICML 2018 report 56.3.} \Textcite{LiuSimonyanYang2019} made the search differentiable. In DARTS each edge of the cell computes a softmax-weighted mixture of the candidate operations; the mixture weights are learned by gradient descent on the validation loss, the network weights on the training loss, and the strongest operations are kept at the end. Second-order DARTS reached $2.76\pm0.09$\% test error on CIFAR-10 (mean and standard deviation of ten trainings of the selected cell) after four GPU-days of search, and 55.7 test perplexity on Penn Treebank after one.

\paragraph{The random arms.} Both papers asked whether the search mattered, which many papers of the period did not \parencite{LindauerHutter2020}. ENAS sampled one architecture uniformly from each of its search spaces and trained it like the searched ones; DARTS ran a random search, drawing 24 cells (eight recurrent cells on Penn Treebank) and keeping the one with the best validation error after a short training.

\begin{finding}[\textcite{PhamEtAl2018}; \textcite{LiuSimonyanYang2019}]\label{fd:controls:claim}
In the ENAS spaces, one uniformly sampled recurrent cell reached test perplexity 81.2 on Penn Treebank, against 56.3 for the ENAS cell, and one pair of sampled convolutional cells reached 6.77\% test error on CIFAR-10, against 3.54\% for the ENAS cells, both trained without cutout. In the DARTS space, the best of 24 random cells, kept for its validation error after 100 epochs, reached $3.29\pm0.15$\% test error, against $2.76\pm0.09$\% for second-order DARTS (means and standard deviations of ten trainings); the best of eight random recurrent cells, kept after 300 epochs, reached perplexity 59.4, against 55.7.
\end{finding}

The DARTS paper drew a measured conclusion: random search was competitive on both tasks, which ``reflects the importance of the search space design'', and DARTS, at comparable or lower cost, improved on it substantially. The field received the stronger half. Weight-sharing searches, many of them variants of DARTS, multiplied, and papers compared themselves in tables of test errors and GPU-days taken from one another. The received claim was that a cheap search finds better architectures than random search does.

\paragraph{What was right.} The cost reduction was real, and it brought architecture search within reach of a university laboratory. Both papers ran a random control and published their code, and DARTS reported the variation of its final trainings. The claim was also reasonable. Searches had produced state-of-the-art classifiers for CIFAR-10 and ImageNet \parencite{ZophEtAl2018,RealEtAl2019}, and in a controlled comparison in the NASNet space, with the same code for every arm and five repeats, evolution and reinforcement learning both compared favorably against random search \parencite{RealEtAl2019}.\sidenote{Each arm of that comparison trained 20{,}000 models for 25 epochs on 450 GPUs, about a week per experiment; the random arm was a random search with the same budget, a sham in the language of \cref{sec:controls:controls}.} What went wrong was not that search is useless, but that the control arms of the cheap searches were too weak to show what the search added.

\section{The experiment: one search against one random draw}\label{sec:controls:experiment}

The comparison behind the received claim is clearest in the DARTS paper's CIFAR-10 table.

\begin{experiment}[Liu, Simonyan and Yang 2019, CIFAR-10]
\field{question}{Does gradient-based search find better cells than random search in the same space?}
\field{treatment}{the search method: second-order DARTS, first-order DARTS, or random search}
\field{space}{the DARTS cells (convolutions, poolings and skip connections on a fixed cell graph)}
\field{selection}{DARTS run four times, random search drawing 24 cells; in each arm the cell with the best validation error after 100 epochs of training from scratch is kept}
\field{held fixed}{the final pipeline: 20 cells, 36 channels, 600 epochs, cutout, path dropout, auxiliary towers}
\field{response}{test error of the kept cell}
\field{replicates}{ten trainings of the kept cell; each arm's whole procedure (4 GPU-days of search) run once}
\field{analysis}{mean and standard deviation of the ten trainings, beside numbers from other papers}
\end{experiment}

The design has a strength that tables of the period often lacked: the random arm is concurrent, with the same space and the same pipeline, so the pipeline cannot explain the gap between 3.29\% and 2.76\%. Its weakness is the unit of replication. The claim is about two procedures, each of which returns a random architecture; the experiment ran each procedure once and trained the returned architecture ten times. The interval $\pm0.09$ describes the trainings of one cell, not the variation of the search, and whether 3.29\% was a typical outcome of a 24-draw random search or an unlucky one, the design cannot say. Nor does it say how good the average architecture of the space is, so the reader cannot split DARTS's 2.76\% into what the space and pipeline give every architecture and what the search adds. The other rows of the table came from other papers, some re-trained under DARTS's setup; comparisons with them are historical. The ENAS controls had one architecture per space.

A reader in 2019 could conclude that DARTS had beaten one draw of random search under one pipeline, but not how large the effect of the search was relative to its own variability, nor how close a randomly sampled cell, trained the same way, would come.

\section{The post-mortem: a fair draw for the random arm}\label{sec:controls:postmortem}

Three studies, published in 2019 and 2020, gave the random arm what it had lacked: more draws, repeated searches, and the same pipeline.

\paragraph{Random search, done properly.} \Textcite{LiTalwalkar2019} treated architecture search as hyperparameter optimization and ran two strong random baselines in the DARTS spaces. Random search with early stopping (the bandit method ASHA, from one to 300 epochs) evaluated more than 700 cells on CIFAR-10 in nine GPU-days and reached $3.03\pm0.13$\% test error (mean and standard deviation of ten trainings), against 3.29\% for the DARTS random baseline; on Penn Treebank it reached perplexity 56.4, against 59.4. Random search with weight sharing, which ranks thousands of random cells with weights trained on uniformly sampled subgraphs, reached 55.5 on Penn Treebank, better than DARTS's 55.7, and $2.85\pm0.08$\% on CIFAR-10, close to DARTS's $2.76\pm0.09$\% and to their own rerun of it, $2.78\pm0.12$\%. They found none of twelve recent search papers exactly reproducible from what was released, and their own method, rerun with new seeds, gave 56.5 instead of 55.5.

\paragraph{The search phase against a placebo.} \Textcite{YuEtAl2020} compared each search with one architecture sampled uniformly from the same space and trained like the searched one, and repeated every arm ten times with different seeds.\sidenote{The paper first circulated in 2019 with Christian Sciuto as first author and is often cited as Sciuto et al.; the published version lists Kaicheng Yu and Christian Sciuto as equal first authors.} On Penn Treebank the random cells reached validation perplexity $60.13\pm0.65$ (mean and standard deviation over the seeds), against $59.88\pm1.92$ for ENAS, $60.61\pm2.54$ for DARTS and $61.99\pm1.95$ for NAO; Welch's $t$-tests gave $p=0.73$, $0.62$ and $0.02$, the last because NAO was worse. On CIFAR-10 the random cells reached $96.44\pm0.19$\% accuracy and the three searches between 96.62\% and 96.86\%, less than half a point better. They traced the weakness to the proxy: in a reduced recurrent space of 32 cells, each trained ten times to obtain its true rank, the Kendall rank correlation between the ranking given by shared weights and the true one averaged $-0.004$ over ten runs, and in NAS-Bench-101 spaces of three to seven nodes it fell from 0.44 to 0.20; samplers trained on true accuracies instead beat random sampling with probability 0.90 (ENAS) and 0.92 (NAO). Their reading was that the spaces were constrained so well that a random architecture already did well.

\paragraph{Eight methods, five datasets, one placebo.} \Textcite{YangEsperancaCarlucci2020} ran eight search methods on five image datasets, eight searches and eight random architectures per method and dataset, all trained with each method's own pipeline. Their score, the relative improvement over the placebo in percent of its accuracy, ranged on CIFAR-10 from $-0.48$ to $0.74$ (DARTS 0.32, ENAS 0.01), and on the other datasets several methods fell below their placebo. The training protocol mattered more: from the plainest pipeline to one with every trick, the same random cells gained more than three points, while the best improvement of a search over random sampling on CIFAR-10 was 0.69 points. Over 214 random cells of the DARTS space, trained with the DARTS pipeline, accuracy was $97.03\pm0.23$\% (mean and standard deviation), from 96.18\% to 97.56\%; a deliberately weaker set of operations moved this distribution by only 0.18 points, and two trainings of the same 32 cells with different seeds ranked them with Kendall $\tau=0.48$.

\paragraph{Worse than the placebo, and better.} In the tabular benchmark NAS-Bench-201, where all 15{,}625 cells have been trained, DARTS with its original hyperparameters converged to a cell made only of skip connections, with 54.30\% test accuracy on CIFAR-10, while random search reached $93.70\pm0.36$\% (mean and standard deviation of 500 runs) and the best cell 94.37\% \parencite{DongYang2020}; \textcite{ZelaEtAl2020} found such degenerate solutions of DARTS in a wide range of spaces. The same literature marks the limits of the verdict: in NAS-Bench-101, where the proxy is full training, regularized evolution and two Bayesian optimizers overtook random search after about 25 evaluated architectures and reached its final performance about five times faster \parencite{YingEtAl2019}. What vanished is the margin of the cheap, weight-sharing searches over a fair random arm in the spaces where they were introduced.

\paragraph{The mechanism.} Four things made a small effect look large. The random arms were underpowered: one architecture per space for ENAS, one random search for DARTS, with selection after a short training that \textcite{LiTalwalkar2019} found unstable. The spaces were narrow by design, so the spread between random architectures was about the size of the noise between seeds. The proxies, shared weights, ranked architectures little better than chance, and a search guided by a proxy that does not rank cannot beat the placebo (\cref{sec:controls:proxy}). And comparisons with published numbers mixed pipelines whose differences were worth more than the search. The same pattern appears wherever new methods are compared with old numbers: given the same tuning and pipeline, the standard LSTM matched or beat newer recurrent cells, the NAS cell among them \parencite{MelisDyerBlunsom2018}, the GAN variants tested did not consistently beat the original \parencite{LucicEtAl2018}, and claimed gains over decade-old metric-learning methods shrank to marginal \parencite{MusgraveBelongieLim2020}. The missing control was the old method given the new method's budget and pipeline (\cref{ch:budgets}).

\begin{verdict}{vanished}
The large margin of the cheap, weight-sharing searches over random architectures did not survive a fair random arm. With the same space, pipeline and enough draws, random architectures matched the searched cells on Penn Treebank and came within a few tenths of a point on CIFAR-10, a remainder of the order of the seed-to-seed variation that ten searches detected for some methods and not for others. The reported gaps came from random arms of one to twenty-four architectures, from proxies that barely ranked, and, in tables of published numbers, from different pipelines. With a proxy that ranks well and a space with real spread, search does beat random search, by a small margin.
\end{verdict}

\section{The method: placebo, sham and the spread of a space}\label{sec:controls:controls}

An experiment learns about a treatment only when its arms differ in nothing else, and the older sciences name the arms that make this true. A \term{negative control}\index{control!negative} is a condition in which the effect cannot occur, and the assay must show none; epidemiologists borrowed the practice from the laboratory to detect bias \parencite{LipsitchEtAl2010}. A \term{positive control}\index{control!positive} is a condition in which the effect is known to occur, and the assay must show it; the guideline on control groups in clinical trials calls the ability of a trial to tell an effective treatment from an ineffective one its \term{assay sensitivity}\index{assay sensitivity} \parencite{ICH2000}. A \term{placebo}\index{control!placebo} is an inert treatment that looks like the real one, and controls for every influence on the course of the disease other than the drug, regression to the mean and expectations included. A \term{sham}\index{control!sham} procedure, in surgery and device trials, is the whole procedure without its active step. The guideline treats the external or \term{historical control}\index{control!historical}, a group treated earlier or elsewhere, separately, because it can rarely guarantee that the groups are comparable.\sidenote{\Textcite{SacksChalmersSmith1982} found six therapies tested both ways: 44 of 56 trials with historical controls (79\%) found the new therapy better, but only 10 of 50 randomized trials (20\%). The treated patients did about equally well in both; the historical control groups did worse.}

A search method has the same anatomy. It consists of a search space $\cA$ with a sampling distribution $\pi$ (uniform, for a finite space); a \term{proxy}\index{proxy}, a cheap and noisy measurement of an architecture's quality such as the validation error after a short training or under shared weights; a budget of $k$ proxy evaluations; a rule that selects the winner; and a final training pipeline that produces the reported number. The active ingredient is the search strategy; everything else is the vehicle.

\begin{definition}[Controls for a search method]\label{def:controls:controls}
For a search method with space $\cA$, distribution $\pi$, proxy, budget $k$, selection rule and final pipeline:
\begin{enumerate}
\item the \emph{placebo arm} draws one architecture from $\pi$ and trains it with the final pipeline (random sampling);
\item the \emph{sham arm}\index{random search} draws $k$ architectures independently from $\pi$, evaluates them with the same proxy, selects with the same rule and trains the winner with the final pipeline (random search with the same budget);
\item a \emph{positive control} is a run in which a gain is known to exist, such as a perfect proxy (an \emph{oracle}) or a planted architecture of known quality, and the experiment must detect it;
\item a \emph{negative control} is a run in which no gain can exist, such as a proxy whose values are permuted among the architectures, and the experiment must report none;
\item a \emph{historical control} is a number produced earlier by another pipeline, usually in another paper.
\end{enumerate}
\end{definition}

The placebo asks whether searching helps at all, and measures the space; the sham asks whether this search strategy helps more than drawing at random with the same proxy and budget.\sidenote{\Textcite{LindauerHutter2020} make the same distinction, between \emph{random sampling}, which returns one random architecture at no search cost, and \emph{random search}, which runs as long as the method it is compared with, and recommend reporting both.} A method that beats its placebo but not its sham owes its gain to the proxy and the budget, which random search exploits equally well. Both arms are concurrent: they share the space, the proxy and the pipeline with the search. A historical control shares none of them. \Cref{ch:ablations} uses the same device for ablations, with a shuffled pruning mask as the sham.

\paragraph{What a historical comparison measures.} Write $\mu_p(a)$ for the true error of architecture $a$ trained by pipeline $p$, the mean of its run distribution over seeds, and $\bar\mu_p$ for the mean of $\mu_p(A)$ over $A\sim\pi$, the expected error of the placebo. A paper that reports its searched architecture $A$ under a new pipeline against a baseline $b$ measured under an old pipeline reports $\mu_{\mathrm{old}}(b)-\mu_{\mathrm{new}}(A)$. Adding and subtracting $\mu_{\mathrm{new}}(b)$ and $\bar\mu_{\mathrm{new}}$ splits it into three parts:
\begin{equation}\label{eq:controls:decomp}
\mu_{\mathrm{old}}(b)-\mu_{\mathrm{new}}(A)=\underbrace{\mu_{\mathrm{old}}(b)-\mu_{\mathrm{new}}(b)}_{\text{pipeline}}+\underbrace{\mu_{\mathrm{new}}(b)-\bar\mu_{\mathrm{new}}}_{\text{baseline against space}}+\underbrace{\bar\mu_{\mathrm{new}}-\mu_{\mathrm{new}}(A)}_{\text{search}} .
\end{equation}
Only the last term belongs to the search, and only a concurrent placebo arm under the new pipeline measures it. The middle term is negative when the baseline is better than the average architecture of the space, as a well-chosen textbook architecture usually is; the first is the co-treatment of \cref{ch:confounding}.

\paragraph{The spread of a space.} Let $Y$ be the test error of one run and $\mu(a)=\E[Y\mid a]$ the \term{true error}\index{true error} of architecture $a$ under the fixed pipeline. For $A\sim\pi$, the mean $\bar\mu$ of $\mu(A)$ is the expected error of the placebo, its standard deviation $\sigma_{\cA}$ is the \term{spread}\index{search space!spread} of the space, and $H=\bar\mu-\mu^\star$, with $\mu^\star$ the smallest true error in the space, is its \term{headroom}\index{search space!headroom}. A search that returns the architecture $A_{\cS}$ has \emph{placebo-adjusted gain} $G=\bar\mu-\E\mu(A_{\cS})$.

\begin{remark}[No search beats the headroom]\label{rem:controls:headroom}
Every search satisfies $G\le H$, because $\mu(A_{\cS})\ge\mu^\star$ for every outcome of the search. In a space whose headroom is a few tenths of a point, no strategy can gain more than that over the placebo, however expensive.
\end{remark}

The spread is measured with noise. If each of $m$ sampled architectures is trained $n$ times and $\bar Y_a$ is its mean, then $\Var\bar Y_a=\sigma_{\cA}^2+\sigma^2/n$, where $\sigma^2$ is the run-to-run variance within an architecture, so the standard deviation of the observed means overstates the spread; the one-way random-effects estimator subtracts the within-architecture mean square divided by $n$ (\cref{thm:toolkit:anova}; the crossed version is in \cref{ch:seeds}). The best observed mean is biased downward like every best-of-$k$ number (\cref{ch:seeds}), so the headroom is estimated honestly only by retraining the winner with fresh seeds.

How many random draws does a perfect search need? The answer does not depend on the size of the space.

\begin{proposition}[Random search reaches a quantile]\label{prop:controls:quantile}
Let $A_1,\dots,A_k$ be independent draws from $\pi$, and let $\cQ\subseteq\cA$ be a set with $\pi(\cQ)=q\in(0,1]$, for instance the best $q$-fraction of the space. Then
\[
  \Prob\bigl(A_i\in\cQ\text{ for some }i\le k\bigr)=1-(1-q)^k\ge1-e^{-kq},
\]
so $k\ge\ln(1/\delta)/q$ draws reach $\cQ$ with probability at least $1-\delta$.
\end{proposition}

\begin{proof}
The events $\{A_i\notin\cQ\}$ are independent, each with probability $1-q$, so all $k$ occur with probability $(1-q)^k$, and the complement is the stated event. The inequality $1-q\le e^{-q}$, which says that the convex function $e^{-q}$ lies above its tangent at $q=0$, gives $(1-q)^k\le e^{-kq}$, and $k\ge\ln(1/\delta)/q$ makes $e^{-kq}\le\delta$. The degenerate case $q=1$ is covered as well: every draw lies in $\cQ$, and the probability is $1-0^k=1$.
\end{proof}

With $q=\delta=0.05$, sixty draws reach the best 5\% of a space with probability $1-0.95^{60}=0.954$, whether the space has a hundred architectures or $10^{18}$.\sidenote{\Textcite{BergstraBengio2012} use the same formula, $1-(1-v/V)^T$ for a target of relative volume $v/V$ and $T$ trials, to compare random search with grids.} But a random search returns a draw from $\cQ$ only when its proxy recognizes it, which is the subject of the next section.

\section{The method: a proxy caps what a search can find}\label{sec:controls:proxy}

A proxy evaluation of an architecture returns a random number $P$, such as the validation error of one short training with a fresh seed. For $A\sim\pi$ and one proxy evaluation of it, $T=\mu(A)$ and $P$ form a random pair $(T,P)$. Random search with budget $k$ draws $k$ independent copies $(T_1,P_1),\dots,(T_k,P_k)$ and returns the architecture with the smallest proxy error. What that does to $\E T$ runs through the \emph{regression function} $m$, defined by $m(P)=\E[T\mid P]$: the expected true error of an architecture given what the proxy says.

\begin{lemma}[Selection acts through the regression]\label{lem:controls:regression}
Let $(T_1,P_1),\dots,(T_k,P_k)$ be independent copies of a pair $(T,P)$ with $\E\abs T<\infty$, and let $m(P)=\E[T\mid P]$. Let the index $J\in\{1,\dots,k\}$ be a function of $P_1,\dots,P_k$ and of a random variable $U$ independent of all the pairs. Then $\E T_J=\E m(P_J)$.
\end{lemma}

\begin{proof}
The plan is to condition on everything the selection sees. Let $\cG$ be the $\sigma$-field generated by $P_1,\dots,P_k$ and $U$; by hypothesis the indicator $\ind\{J=i\}$ is $\cG$-measurable. The pair $(T_i,P_i)$ is independent of the other proxies and of $U$, so conditioning on them adds nothing to conditioning on $P_i$, and $\E[T_i\mid\cG]=m(P_i)$ almost surely. In the display below, the first equality is the tower property with the $\cG$-measurable indicator taken out of the conditional expectation, legitimate because $\abs{T_i\ind\{J=i\}}\le\abs{T_i}$ is integrable, and the second is the identity $\E[T_i\mid\cG]=m(P_i)$:
\[
  \E\bigl[T_i\ind\{J=i\}\bigr]=\E\bigl[\ind\{J=i\}\,\E[T_i\mid\cG]\bigr]=\E\bigl[\ind\{J=i\}\,m(P_i)\bigr].
\]
Summing over $i$ gives $\E T_J=\sum_i\E[T_i\ind\{J=i\}]=\E m(P_J)$.
\end{proof}

The lemma is the theory of the negative control: if the proxy carries no information, so that $m$ is constant, then $\E T_J=\bar\mu$ for every rule that looks only at proxy values, however clever, and the search is the placebo. With a Gaussian proxy the gain can be computed. Recall from \cref{ch:seeds} that $e_k$ is the expected maximum of $k$ independent standard normal variables.

\begin{theorem}[Selection on a Gaussian proxy]\label{thm:controls:gauss}
Let $(T,P)$ be bivariate normal with means $\mu_T,\mu_P$, standard deviations $\sigma_T,\sigma_P>0$ and correlation $\rho$, let $(T_1,P_1),\dots,(T_k,P_k)$ be independent copies, and let $J$ be the index of the smallest $P_i$. Then
\begin{equation}\label{eq:controls:gain}
  \E T_J=\mu_T-\rho\,\sigma_T\,e_k\qquad\text{and}\qquad\E P_J=\mu_P-\sigma_P\,e_k .
\end{equation}
In the additive model $P=T+\xi$, with $\xi\sim\cN(0,s^2)$ independent of $T$, the correlation is $\rho=\sigma_T/\sqrt{\sigma_T^2+s^2}$.
\end{theorem}

\begin{proof}
The plan is to reduce the statement to the best-of-$k$ computation of \cref{ch:seeds}, which is the Gaussian case of \cref{lem:controls:regression}. The $P_i$ have a continuous distribution, so ties have probability zero and $J$ is almost surely unique.

\emph{The selected architecture.} The index $J$ of the smallest $P_i$ is the index of the largest $D_i=-P_i$, and the pairs $(D_i,T_i)$ are independent and bivariate normal with standard deviations $\sigma_P,\sigma_T$ and correlation $-\rho$. Part~(ii) of \cref{prop:seeds:bestof}, whose proof computes the regression function $m$ and applies the reasoning of \cref{lem:controls:regression}, gives $\E T_J=\mu_T+(-\rho)\sigma_Te_k=\mu_T-\rho\sigma_Te_k$.

\emph{The selected proxy.} Part~(i) of the same proposition, applied to the $P_i$, gives $\E P_J=\E\min_iP_i=\mu_P-\sigma_Pe_k$.

\emph{The additive model.} Independence of $T$ and $\xi$ gives $\Cov(T,P)=\sigma_T^2$ and $\sigma_P^2=\sigma_T^2+s^2$, so $\rho=\sigma_T^2/(\sigma_T\sigma_P)=\sigma_T/\sqrt{\sigma_T^2+s^2}$.
\end{proof}

With $T=\mu(A)$, so that $\mu_T=\bar\mu$ and $\sigma_T=\sigma_{\cA}$, random search with a Gaussian proxy has placebo-adjusted gain $G=\rho\,\sigma_{\cA}\,e_k$. There are three ways to be no better than the placebo: a narrow space ($\sigma_{\cA}$ small), a small budget ($e_k$ grows only like $\sqrt{2\ln k}$), or a proxy that does not rank ($\rho$ near zero).

\begin{corollary}[Efficiency and the winner's curse]\label{cor:controls:curse}
Under the hypotheses of \cref{thm:controls:gauss}:
\begin{enumerate}
\item the placebo-adjusted gain of the proxy-guided search is the fraction $\rho$ of the gain $\sigma_Te_k$ of the same search with a perfect proxy;
\item in the additive model, the gain that the search reports on its own proxy, $\mu_P-\E P_J$, is the true gain divided by $\rho^2$.
\end{enumerate}
\end{corollary}

\begin{proof}
A perfect proxy selects the smallest $T_i$, and the computation of the minimum in the proof of \cref{thm:controls:gauss}, applied to the $T_i$, gives $\E\min_iT_i=\mu_T-\sigma_Te_k$; dividing $\rho\sigma_Te_k$ by $\sigma_Te_k$ proves (i). In the additive model $\E\xi=0$ gives $\mu_P=\mu_T$, and $\sigma_P=\sigma_T/\rho$, so the reported gain is $\sigma_Pe_k=\sigma_Te_k/\rho$ while the true gain is $\rho\sigma_Te_k$; their ratio is $1/\rho^2$.
\end{proof}

Part (ii) is the \term{winner's curse}\index{winner's curse} of a search: the winner's proxy score is the best of $k$ noisy numbers, and its noise was selected along with its quality. With $\rho=0.5$ a search reports four times the gain it achieved. Animal breeders know \cref{thm:controls:gauss} in the additive case as the breeder's equation $R=h^2S$: the response to selection is the heritability times the selection differential.\sidenote{The equation goes back to Jay Lush's \emph{Animal Breeding Plans} of 1937; \textcite{FalconerMackay1996} derive it. Here $S=\sigma_Pe_k$, $h^2=\rho^2$, and $R=\rho\sigma_Te_k$.}

Search papers report the quality of a proxy as a rank correlation, Kendall's $\tau$\index{Kendall's tau@Kendall's $\tau$} \parencite{Kendall1938}: the probability that two architectures are ordered the same way by the proxy and by the truth, minus the probability that they are ordered oppositely (\cref{thm:toolkit:kendall}). For Gaussian pairs it determines $\rho$.

\begin{lemma}[Kendall's $\tau$ of a Gaussian pair]\label{lem:controls:kendall}
Let $(X,Y)$ be bivariate normal with positive variances and correlation $\rho$, and let $(X_1,Y_1),(X_2,Y_2)$ be independent copies. Then
\[
  \tau=\Prob\bigl((X_1-X_2)(Y_1-Y_2)>0\bigr)-\Prob\bigl((X_1-X_2)(Y_1-Y_2)<0\bigr)=\frac2\pi\arcsin\rho .
\]
\end{lemma}

\begin{proof}
The plan is to reduce to a pair of standard normals and read the answer off an angle. The differences $U=(X_1-X_2)/\sqrt{2\Var X}$ and $V=(Y_1-Y_2)/\sqrt{2\Var Y}$ are jointly normal with unit variances and, because the copies are independent, correlation $\rho$; they have the signs of $X_1-X_2$ and $Y_1-Y_2$. If $\rho=\pm1$, then $V=\pm U$ almost surely and $\tau=\pm1=\frac2\pi\arcsin(\pm1)$. Let $\abs\rho<1$ and $\varphi=\arcsin\rho\in(-\pi/2,\pi/2)$, so that $\sqrt{1-\rho^2}=\cos\varphi$. The pair $(U,V)$ has the law of $(Z_1,Z_1\sin\varphi+Z_2\cos\varphi)$ with independent standard normals $Z_1,Z_2$; write $(Z_1,Z_2)=\Lambda(\cos\Theta,\sin\Theta)$ in polar coordinates, with $\Theta$ uniform on $[0,2\pi)$ and $\Lambda>0$. Then $U=\Lambda\cos\Theta$ and $V=\Lambda\sin(\Theta+\varphi)$. Both are positive when $\Theta\in(-\pi/2,\pi/2)$ and $\Theta+\varphi\in(0,\pi)$ modulo $2\pi$, that is, on the arc $(-\varphi,\pi/2)$ of length $\pi/2+\varphi$, and both are negative on the opposite arc, of the same length. Hence $\Prob(UV>0)=\frac12+\varphi/\pi$, the event $UV=0$ has probability zero, and $\tau=2\Prob(UV>0)-1=2\varphi/\pi$.
\end{proof}

Inverting, $\rho=\sin(\pi\tau/2)$: $\tau=0.2$ corresponds to $\rho=0.31$, and $\tau=0.5$ to $\rho=0.71$.

\begin{example}[Reading published correlations]\label{exa:controls:reading}
The weight-sharing ranking of \textcite{YuEtAl2020} in their 32-cell recurrent space had $\tau=-0.004$, so $\rho\approx0$, and \cref{thm:controls:gauss} predicts no gain over the placebo, which is what they found. In the seven-node NAS-Bench-101 space $\tau=0.195$ gives $\rho=0.30$: the search captures 30\% of what a perfect search with the same draws would. Two single trainings of the same 32 DARTS cells had $\tau=0.48$ \parencite{YangEsperancaCarlucci2020}, a correlation $\rho_2=\sin(0.24\pi)=0.68$. If each training is the true quality plus independent noise, then $\rho_2=\sigma_{\cA}^2/(\sigma_{\cA}^2+\sigma^2)$, and their standard deviation of 0.23 points across cells splits into a spread $\sigma_{\cA}\approx0.19$ and a seed noise $\sigma\approx0.13$ points. A perfect search of 24 random cells, the size of the DARTS random arm, would gain $\sigma_{\cA}e_{24}\approx0.19\times1.95=0.37$ points over the average cell; a proxy with $\tau=0.2$, about 0.11, less than the noise of one training. These are orders of magnitude, resting on a Gaussian model and another group's implementation.
\end{example}

The theorem assumes that the regression of the truth on the proxy is a straight line all the way into the tail where the winners come from. Real proxies have a second, fixed error, which bends it.

\begin{heuristic}[A biased proxy reaches a plateau]
Let a proxy evaluation of $a$ be $\mu(a)+b(a)+\varepsilon$, with a bias $b(a)$ that belongs to the architecture (a short training favors fast learners) and noise $\varepsilon$ that is fresh at each evaluation. In a finite space, as the budget grows the draws cover the space and the winner tends to the minimizer of $\mu+b$, perturbed by the selected noise; its expected true error tends to a limit set by how well $\mu+b$ ranks the architectures, not by $k$, short of the headroom. A strategy that concentrates its evaluations on the proxy's favorites, as evolution does, reaches this plateau sooner, and when the bias is large a better optimizer of the proxy can end with a worse architecture than random search.
\end{heuristic}

\begin{remark}[How many searches]\label{rem:controls:power}
The unit of a search experiment is the search. With $n$ searches per arm and $r$ trainings of each returned architecture, an arm's mean has variance $(\sigma_S^2+\sigma^2/r)/n$, where $\sigma_S$ is the standard deviation of the true error of what the search returns ($\sigma_S=\sigma_{\cA}$ for the placebo). The sample-size formula of \cref{thm:toolkit:power} gives $n\approx2(z_{0.975}+z_{0.8})^2(\sigma_S^2+\sigma^2/r)/G^2$ per arm for power $0.8$ at level $0.05$. With $\sigma_S=0.2$, $\sigma=0.13$ and $r=1$ (\cref{exa:controls:reading}), a gain of $G=0.3$ points needs about ten searches per arm, and a gain of 0.2 points about 22. \end{remark}

\section{In the laboratory: a tabular benchmark of 192 networks}\label{sec:controls:lab}

Architecture search on the digits data is cheap enough to do exhaustively: we train every architecture of a small space several times, store the results in a \term{tabular benchmark}\index{tabular benchmark}, and run thousands of searches as table lookups, as NAS-Bench-101 does for CIFAR-10 \parencite{YingEtAl2019}. The data are labkit's split 0 of the scikit-learn digits; its 1258 training images are divided once into 958 for training and 300 for validation, and the 539 test images, each worth 0.19 points of test error, are used only to define the truth.

\begin{lab}{The spread of a search space}{05-controls-space}\label{lab:controls:space}
\labpart{Setup} Every MLP with depth 1 to 4, width 16, 32, 64 or 128, activation ReLU, tanh or GELU, dropout 0 or 0.25, and batch normalization on or off: 192 architectures. One fixed pipeline: Adam with learning rate $10^{-3}$, batch 64, 30 epochs at a constant learning rate, on the 958 training images. Each architecture is trained six times (seeds 2000--2005) and scored on the test images; the mean of the six is its true error. 1152 runs, three to four minutes.
\labpart{Predict first} How far is the average architecture from the best? Is the spread between architectures larger or smaller than the spread between two seeds of one architecture? Answer for the full space and for the 96 architectures of width 64 or 128.
\labpart{Look for} The shape of the ranked curve, the overlap of the intervals near the top, and where the wide architectures sit.
\end{lab}

\begin{figure}[t]
\begin{fullwidth}
\centering
\begin{tikzpicture}
\begin{axis}[lplot,width=\linewidth,height=0.3\linewidth,ymode=log,log basis y=10,
  xmin=-2,xmax=195,ymin=1,ymax=45,ytick={1,2,3,5,10,20,40},yticklabels={1,2,3,5,10,20,40},
  xlabel={architecture, ranked by mean test error},ylabel={test error (\%)},
  legend pos=north west,legend columns=1]
\addplot[run,mark size=0.85pt] table[x=rank,y expr=100*\thisrow{err}] {figures/data/05-controls-space-runs.dat};
\addlegendentry{one run}
\addplot[mean,mark size=1.15pt,ebars] table[x=rank,y expr=100*\thisrow{meanwide},y error expr=100*\thisrow{hw}] {figures/data/05-controls-space.dat};
\addlegendentry{mean, width 16 or 32}
\addplot[only marks,mark=square*,mark size=1.05pt,color=rubric,mark options={fill=rubric,draw=white,line width=0.15pt},
  error bars/y dir=both,error bars/y explicit,error bars/error bar style={rubric,line width=0.45pt},
  error bars/error mark options={rubric,rotate=90,mark size=1.2pt,line width=0.45pt}]
  table[x=rank,y expr=100*\thisrow{meancur},y error expr=100*\thisrow{hw}] {figures/data/05-controls-space.dat};
\addlegendentry{mean, width 64 or 128}
\addplot[reference] table[x=x,y expr=100*\thisrow{full}] {figures/data/05-controls-space-ref.dat};
\addplot[reference,rubric] table[x=x,y expr=100*\thisrow{curated}] {figures/data/05-controls-space-ref.dat};
\end{axis}
\end{tikzpicture}
\caption{The spread of the search space of \cref{lab:controls:space}: 192 MLP architectures for the digits data, ranked by mean test error over six runs (seeds 2000--2005, Adam, 30 epochs). Hollow marks: the six runs of each architecture. Filled marks: means with 95\% $t$-intervals, vermilion squares for the 96 curated architectures (width 64 or 128) and ink circles for the rest. Dashed lines: the expected error of a randomly sampled architecture, 4.64\% in the full space (grey) and 3.14\% in the curated space (vermilion). Logarithmic vertical axis.}\label{fig:controls:space}
\end{fullwidth}
\end{figure}

\paragraph{What the space looks like.} \Cref{fig:controls:space} shows two spaces in one. The full space has a long tail: its mean true error is 4.64\%, the standard deviation of its means 3.18 points, and its worst architecture, four layers of 16 ReLU units with dropout, errs on 27.33\% of the test images. Width decides most of it (mean error 7.99\% at width 16, 2.91\% at width 128), then batch normalization and dropout. The best architecture, two layers of 128 GELU units with batch normalization, reaches \pmval{1.89}{0.38}\% (mean and 95\% $t$-interval over six runs), so the headroom of the full space is 2.76 points. The 96 wide architectures form the kind of space a designer writes down first: its mean is 3.14\%, its headroom 1.26 points, and the random-effects estimate of its spread, $\sigma_{\cA}=0.62$ points, is about 1.3 times the seed noise within an architecture, $\sigma=0.47$ points; in the DARTS space the ratio was about $0.19/0.13\approx1.5$ (\cref{exa:controls:reading}). Two single runs rank the curated architectures with Kendall $\tau=0.52$ on average (0.39 to 0.66 over the fifteen pairs of seeds), against 0.48 in the DARTS space, and eleven of the 95 other curated architectures cannot be told apart from the best by a Welch test at level 0.05.

\begin{lab}{A search, its sham and its controls}{05-controls-search}\label{lab:controls:search}
\labpart{Setup} Four more runs of each architecture (seeds 1000--1003) record the validation error after every epoch: the proxy of fidelity $f$ is the error after $f$ epochs, and the main proxy is $f=2$. An evaluation draws one of the four runs at random, as NAS-Bench-101 does. Five arms, each repeated 2000 times with $k=1,2,4,\dots,64$ evaluations: the placebo, random search (the sham), regularized evolution (population 10, tournament 3, one-attribute mutations), the oracle (best of $k$ draws by true error, a positive control) and random search with a shuffled proxy (a negative control). A search is scored by the true error of what it returns. 768 runs, about three minutes.
\labpart{Predict first} In the curated space the two-epoch proxy has Kendall $\tau=0.32$ with the truth. What fraction of the oracle's gain will random search capture with 16 evaluations? Will evolution beat random search?
\labpart{Look for} Where each curve flattens, the gap between random search and the oracle, and what the shuffled search reports about itself.
\end{lab}

\begin{figure}[t]
\begin{fullwidth}
\centering
\begin{tikzpicture}
\begin{groupplot}[group style={group size=2 by 1,horizontal sep=15mm},lplot,
  width=0.5\linewidth,height=0.32\linewidth,xmode=log,log basis x=2,xmin=0.75,xmax=88,
  xtick={1,2,4,8,16,32,64},xticklabels={1,2,4,8,16,32,64},minor tick num=0,
  xlabel={proxy evaluations $k$},unbounded coords=jump]
\nextgroupplot[ylabel={true error of returned architecture (\%)},ymin=1.8,ymax=5.2,
  title={(a) full space, 192 architectures}]
\addplot[run,restrict expr to domain={\thisrow{space}}{0:0}] table[x=x,y=err] {figures/data/05-controls-search-runs.dat};
\addplot[reference] table[x=k,y=mean_full] {figures/data/05-controls-search-curves.dat};
\addplot[black!55,line width=0.6pt,densely dotted] table[x=k,y=best_full] {figures/data/05-controls-search-curves.dat};
\addplot[extrapolation] table[x=k,y=pred_full] {figures/data/05-controls-search-curves.dat};
\addplot[black!65,line width=0.9pt,mark=triangle*,mark size=1.8pt,mark options={fill=black!65,draw=white,line width=0.2pt}] table[x=k,y=oracle_full] {figures/data/05-controls-search-curves.dat};
\addplot[armB,error bars/y dir=both,error bars/y explicit,error bars/error bar style={rubric,line width=0.5pt}] table[x=k,y=evolution_full,y error=evolution_full_hw] {figures/data/05-controls-search-curves.dat};
\addplot[armA,ebars] table[x=k,y=random_search_full,y error=random_search_full_hw] {figures/data/05-controls-search-curves.dat};
\nextgroupplot[ymin=1.8,ymax=3.4,title={(b) curated space, width 64 or 128},
  legend to name=leg:controls:search,legend columns=4,
  legend style={draw=none,fill=none,font=\footnotesize,/tikz/every even column/.append style={column sep=8pt}}]
\addplot[run,restrict expr to domain={\thisrow{space}}{1:1}] table[x=x,y=err] {figures/data/05-controls-search-runs.dat};
\addlegendentry{one random search}
\addplot[reference] table[x=k,y=mean_curated] {figures/data/05-controls-search-curves.dat};
\addlegendentry{placebo}
\addplot[black!55,line width=0.6pt,densely dotted] table[x=k,y=best_curated] {figures/data/05-controls-search-curves.dat};
\addlegendentry{best architecture}
\addplot[extrapolation] table[x=k,y=pred_curated] {figures/data/05-controls-search-curves.dat};
\addlegendentry{Gaussian prediction}
\addplot[black!65,line width=0.9pt,mark=triangle*,mark size=1.8pt,mark options={fill=black!65,draw=white,line width=0.2pt}] table[x=k,y=oracle_curated] {figures/data/05-controls-search-curves.dat};
\addlegendentry{oracle}
\addplot[armB,error bars/y dir=both,error bars/y explicit,error bars/error bar style={rubric,line width=0.5pt}] table[x=k,y=evolution_curated,y error=evolution_curated_hw] {figures/data/05-controls-search-curves.dat};
\addlegendentry{evolution}
\addplot[armA,ebars] table[x=k,y=random_search_curated,y error=random_search_curated_hw] {figures/data/05-controls-search-curves.dat};
\addlegendentry{random search}
\end{groupplot}
\end{tikzpicture}\par\smallskip
\pgfplotslegendfromname{leg:controls:search}
\caption{Searches on the benchmark of \cref{lab:controls:search}, guided by the two-epoch proxy (four runs per architecture, seeds 1000--1003). Vertical axis: the true error (mean of six test runs) of the architecture a search returns. Lines with marks: means over 2000 searches with 95\% intervals (most smaller than the marks), for random search (ink), regularized evolution (vermilion) and the oracle that selects by true error (grey). Hollow marks: 40 single random searches at each $k$, spread sideways. Dashed grey: the placebo; dotted: the best architecture; dashed vermilion: \cref{thm:controls:gauss} for random search, from the space's mean and standard deviation and the proxy's correlation with the truth ($\rho=0.66$ in the full space, 0.45 in the curated one). The shuffled-proxy search (not drawn) stays at the placebo level.}\label{fig:controls:search}
\end{fullwidth}
\end{figure}

\paragraph{The proxy.} In the curated space the two-epoch proxy has Kendall $\tau=0.32$ with the truth (Pearson $\rho=0.45$); one epoch gives 0.20, ten epochs 0.46, and the thirty-epoch validation error 0.56, close to the 0.52 at which two full test runs agree. Its bias shows in its favorite: averaged over its four runs, the two-epoch proxy prefers four layers of 128 tanh units with batch normalization, a fast learner that is 33rd of 96 by true error (2.81\%), while the thirty-epoch validation error prefers three layers of 128 GELU units with batch normalization, which is second (2.04\%). In the full space the proxies rank better ($\tau=0.57$ at two epochs), because the narrow networks of the tail are bad early and late.

\paragraph{The search against its controls.} In the full space (\cref{fig:controls:search}) every search with an informative proxy beats the placebo by about two points, because a single random draw lands too often in the tail: with 16 evaluations random search returns architectures with true error \pmval{2.73}{0.02}\% (mean and 95\% interval over 2000 searches) against the placebo's 4.64\%. In the curated space the gain is about half a point, \pmval{2.64}{0.02}\% against 3.14\%, the size of one run's seed noise, and single searches vary around that mean with a standard deviation of 0.47 points. In both spaces random search levels off near 2.6\%, while the oracle goes on to 2.13\% at $k=16$ and 1.98\% at $k=64$ in the curated space (intervals below 0.01 points): the plateau of the biased-proxy heuristic. Evolution, the more capable optimizer, does not beat the sham. Its architectures are worse than those of random search by \pmval{0.03}{0.03}, \pmval{0.06}{0.02} and \pmval{0.09}{0.02} points at $k=16$, 32 and 64 in the curated space (\pmval{0.02}{0.03}, \pmval{0.06}{0.03} and \pmval{0.09}{0.02} in the full space): it finds the proxy's favorites faster, and they are fast learners. Given a proxy that ranks better, evolution wins by a little: with the ten- and thirty-epoch proxies it is ahead in seven of the eight comparisons we ran, by 0.02 to 0.07 points (with the thirty-epoch proxy and 64 evaluations, 2.14\% against 2.21\% in the curated space, a difference of \pmval{0.07}{0.01} points).

\paragraph{The theorem.} In the curated space \cref{thm:controls:gauss} predicts random search well while the budget is small: 2.98\%, 2.84\%, 2.73\% and 2.63\% for $k=2$, 4, 8 and 16, against 2.96\%, 2.85\%, 2.71\% and 2.64\% observed. Beyond that the prediction keeps falling (2.46\% at $k=64$) and the observations do not: the smallest proxy values belong to the fast learners, and the straight regression line fails in the tail. In the full space the Gaussian model fails from the start, predicting 1.68\% at $k=8$, below the best architecture, against 2.82\% observed: the long tail inflates the standard deviation to 3.18 points while the headroom is 2.76, and \cref{rem:controls:headroom} is the bound that holds. \Cref{fig:controls:gain} varies the proxy instead of the budget. In the curated space the gain of random search with 16 evaluations, divided by the perfect-proxy gain $\sigma_{\cA}e_{16}$, follows the line $\rho$ from the shuffled proxy ($\rho=0$, gain \pmval{0.00}{0.03} points, a 95\% interval over 2000 searches) through the one-epoch proxy ($\rho=0.28$, \pmval{0.39}{0.02} points against a prediction of 0.32) to full training ($\rho=0.71$, \pmval{0.82}{0.01} against 0.81); in the full space the points lie at 0.33 to 0.42 of $\sigma_{\cA}e_{16}$ whatever the proxy.

\begin{figure}[t]
\centering
\begin{tikzpicture}
\begin{axis}[lplot,width=0.82\linewidth,height=0.4\linewidth,xmin=-0.05,xmax=1.05,ymin=-0.08,ymax=1.08,
  xlabel={correlation $\rho$ of one proxy evaluation with the true error},
  ylabel={gain$/(\sigma_{\cA}e_{16})$},legend pos=north west]
\addplot[fit,domain=0:1,samples=2] {x};
\addlegendentry{prediction $\rho$}
\addplot[armB,only marks,error bars/y dir=both,error bars/y explicit,error bars/error bar style={rubric,line width=0.5pt}] table[x=rho_curated,y=ng_curated,y error=ng_curated_hw] {figures/data/05-controls-search-fidelity.dat};
\addlegendentry{curated space}
\addplot[armA,only marks,ebars] table[x=rho_full,y=ng_full,y error=ng_full_hw] {figures/data/05-controls-search-fidelity.dat};
\addlegendentry{full space}
\end{axis}
\end{tikzpicture}
\caption{The proxy theorem across proxies. Each point is random search with $k=16$ evaluations on the benchmark of \cref{lab:controls:search} (2000 searches, 95\% interval): its gain over the placebo divided by $\sigma_{\cA}e_{16}$ (1.14 points in the curated space, 5.61 in the full space), against the Pearson correlation of one proxy evaluation with the true error. From left to right: the shuffled proxy, the validation error after 1, 2, 3, 5, 10, 20 and 30 epochs, and the true error itself (the oracle). The line is \cref{thm:controls:gauss}.}\label{fig:controls:gain}
\end{figure}

\paragraph{What the search says about itself.} With the thirty-epoch proxy, whose values are validation errors on the scale of the truth, random search with eight evaluations in the curated space reports a gain of 1.64 points over the average proxy score and achieves 0.68: a ratio of 2.39, somewhat more than the $1/\rho^2=2.01$ of \cref{cor:controls:curse}. The negative control needs no model: with the proxy values shuffled among the architectures, random search with 16 evaluations achieves \pmval{-0.03}{0.03} points (95\% interval) and reports 2.09 (in the full space, \pmval{-0.04}{0.14} and 3.44). A search's report is no evidence until it is measured against a control.

\begin{lab}{A historical control}{05-controls-pipeline}\label{lab:controls:pipeline}
\labpart{Setup} A $2\times2$ factorial: the textbook MLP (two layers of 64 ReLU units) and the architecture found in \cref{lab:controls:search} (the curated architecture with the lowest validation error after full training), crossed with the old pipeline of \cref{lab:controls:space} and a new one (Adam with learning rate $2\times10^{-3}$, cosine schedule, 60 epochs, label smoothing 0.1); ten runs per cell, seeds 3000--3009. A placebo arm: 16 random curated architectures, two runs each, under both pipelines. 104 runs, under a minute.
\labpart{Predict first} A paper reports the found architecture with the new pipeline and copies the textbook MLP's number from an older paper that used the old one. How much of the reported improvement belongs to the search?
\labpart{Look for} The pipeline effect on each architecture, and the placebo arm under the new pipeline.
\end{lab}

\begin{table}[t]
\centering
\caption{\Cref{lab:controls:pipeline}: test error (\%), mean and 95\% $t$-interval; ten runs per architecture and pipeline. The placebo row averages 16 random curated architectures, two runs each, with an interval over the 16.}\label{tab:controls:pipeline}
\begin{tabular}{@{}lcc@{}}
\toprule
 & old pipeline & new pipeline\\
\midrule
textbook MLP ($2\times64$, ReLU) & \pmval{3.19}{0.23} & \pmval{2.00}{0.22}\\
found ($3\times128$, GELU, batch norm.) & \pmval{2.04}{0.34} & \pmval{1.69}{0.22}\\
placebo (random curated) & \pmval{2.97}{0.29} & \pmval{2.32}{0.29}\\
\bottomrule
\end{tabular}
\end{table}

\paragraph{The historical comparison.} The paper of the lab sheet would report an improvement of \pmval{1.50}{0.30} points (\cref{tab:controls:pipeline}; 95\% Welch intervals for all differences here). The new pipeline alone improves the textbook MLP by \pmval{1.19}{0.30} points; under the new pipeline the textbook MLP is better than the average random architecture by $2.32-2.00=0.32$ points; and the found architecture beats the concurrent placebo by \pmval{0.63}{0.35} points. These are the three terms of \eqref{eq:controls:decomp}, $1.19-0.32+0.63=1.50$: the search's own share is less than half of what the historical comparison credits to it. Pipeline and architecture also interact: the new pipeline helps the textbook MLP by 1.19 points but the found architecture by only \pmval{0.35}{0.38}, so the contribution of the architecture depends on the pipeline at which it is measured (\pmval{1.15}{0.38} points under the old one, \pmval{0.32}{0.29} under the new), the order dependence of \cref{ch:confounding}. In miniature this is the finding of \textcite{YangEsperancaCarlucci2020}.

\paragraph{What the miniature reproduces and what it does not.} The benchmark reproduces the mechanism: a proxy that ranks poorly holds every search to a plateau, a better optimizer of a biased proxy does not help, a search's report overstates its gain, and a historical control mixes the pipeline into the search. It does not reproduce the near-zero margin over the placebo: our curated space is relatively wider than the DARTS space (headroom 40\% of the mean error, against about 18\% for the 214 cells of \textcite{YangEsperancaCarlucci2020}), and our short-training proxy ranks better than shared weights did, so random search captures about half of the oracle's gain. Weaker proxies move the gain toward the placebo along the line of \cref{fig:controls:gain}; we did not implement weight sharing.

\section{Protocol}\label{sec:controls:protocol}

\begin{protocol}{Controls for a search method}\label{prot:controls:search}
\item Fix the search space, its sampling distribution, the proxy, the budget and the final pipeline before searching, and use them in every arm.
\item Run a placebo arm: train at least ten architectures drawn from the space with the final pipeline, and report their mean, standard deviation and best, the spread and headroom of the space.
\item Run a sham arm: random search with the same proxy, budget and selection rule.
\item Repeat every search with independent seeds, as often as \cref{rem:controls:power} requires for the gain you claim, and report the distribution of what the searches return, not the best search.
\item Score every returned architecture by retraining it with fresh seeds on test data never used for selection; never report the winner's proxy score as its quality.
\item Measure the proxy: report Kendall's $\tau$ between the proxy and full training on a sample of architectures, beside the $\tau$ between two seeds of full training.
\item Run a positive control (an oracle or a planted better architecture, which must be detected) and a negative control (a shuffled proxy, which must show no gain).
\item Use published numbers only as context; claim gains only against concurrent arms, with and without training tricks.
\item Release the space, the pipeline, the seeds and every sampled architecture with its scores.
\end{protocol}

The protocol measures a search within a space, not the space: a search with no placebo-adjusted gain in an excellent space may still be the right product, and a large gain over a poor placebo is worth little, since a designer can always widen a space with bad options, as our full space shows. Transfer is outside its reach: a cell searched on CIFAR-10 and deployed on ImageNet needs controls at the scale of deployment. And it is expensive when a search is: at a GPU-day per search, ten searches in each of three arms cost a month. Then the comparison of strategies belongs on a tabular or surrogate benchmark, and the full experiment can be reduced to the placebo arm, which has no search cost, the proxy correlation, and enough trainings of the final architecture to state its error with an interval.

\begin{exercises}
\item \consolidation (Expected rank of the best of $k$.) In a space of $S$ architectures ranked by true error, let $\nu_k$ be the rank of the best of $k$ uniform draws with replacement. (a) Show that $\Prob(\nu_k\ge r)=\bigl((S-r+1)/S\bigr)^k$ and $\E\nu_k=\sum_{j=1}^S(j/S)^k\le1+S/(k+1)$, and compute $\E\nu_k$ for $S=96$ and $k=16$. (b) Compare with the true rank of what random search returned in \cref{lab:controls:search}, and explain the difference with \cref{lem:controls:regression}.\label{ex:controls:rank}

\item \consolidation (What survives without normality.) (a) Show with \cref{lem:controls:regression} that if $\E[T\mid P]=\alpha+\beta P$, then every proxy-based selection among independent pairs has $\E T_J=\E T+\beta(\E P_J-\E P)$. (b) Let $P\sim\cN(0,1)$ and $T=P^2$. Show that $\Corr(T,P)=0$, that for $k\ge3$ keeping the smallest $P_i$ gives $\E T_J>\E T$ (equality for $k=2$), and that keeping the smallest $\abs{P_i}$ gives $\E T_J<\E T$. What does this say about summarizing a proxy by one correlation?\label{ex:controls:nonnormal}

\item \consolidation (The winner's curse in a null space.) In a space whose architectures are all equally good ($\sigma_T=0$), a search with the additive proxy $P=T+\xi$, $\xi\sim\cN(0,s^2)$, keeps the smallest of $k$ proxy values. Show that it reports a gain of $s\,e_k$ and achieves none, compute the report for $s=0.47$ points and $k=16$, and relate it to the best-of-$k$ reporting of \cref{ch:seeds}.\label{ex:controls:null}

\item \consolidation\hard (Two seeds and one truth.) In the additive model $P_j=T+\xi_j$ with independent $\xi_1,\xi_2\sim\cN(0,s^2)$, show that $\Corr(P_1,P_2)=\Corr(T,P_1)^2$, and check it on the curated space, where two seeds give $\tau=0.52$ and the thirty-epoch proxy has $\rho=0.71$. \hint{the proxy is a validation error and the truth a six-run test mean, so expect agreement in order of magnitude only.}\label{ex:controls:twoseeds}

\item \labwork (Make the proxy better or worse.) (a) In \script{05-controls-search.py} set \texttt{FID\_MAIN} to 1, 5 and 30 and record the gains of random search and of evolution over the placebo at $k=64$ in the curated space. At which fidelity does evolution start to win? Explain with the biased-proxy heuristic. (b) At equal compute, compare 60 two-epoch evaluations with 4 thirty-epoch ones in both spaces. Why does the answer differ between the spaces?\label{ex:controls:fidelity}

\item \labwork (A weaker proxy.) Add a proxy computed from one epoch on half the training images to \cref{fig:controls:gain}. Does it fall on the line?\label{ex:controls:weaker}

\item \labwork (Widen the space.) Add width 8 and dropout 0.5 to the space of \cref{lab:controls:space}. Predict, then measure, how the placebo, random search and the oracle move. Which control does the new tail affect, and which arm of a comparison does it flatter?\label{ex:controls:widen}

\item \design (Controls for data pruning.) A method keeps 30\% of a training set, and its authors claim that training on the kept subset beats training on a random 30\%. Design the placebo, the sham when there is one, a positive control and a negative control, with the number of runs per arm, for a budget of 200 training runs.\label{ex:controls:pruning}

\item \design (A search experiment under a budget.) A search costs one GPU-day, a final training half a GPU-day, and you have 40 GPU-days. You expect a gain of 0.3 points over random sampling, with single trainings varying by 0.15 points and returned architectures by 0.2. Allocate the budget among searches, placebo draws and final trainings, and state the power of your comparison.\label{ex:controls:budget}

\item \critique (A table without a random row.) A paper reports its search method's test error ($2.61\pm0.05$\% over five trainings of its best cell) beside the published errors and search costs of six earlier methods, with no random architectures, and says the method ``outperforms all baselines at a fraction of the cost''. List what the table cannot show, and the fewest extra runs that would make the claim testable. (The description is invented.)\label{ex:controls:critique1}

\item \critique (A historical baseline.) A paper reports 1.8\% error for a new architecture with its pipeline (cosine schedule, 300 epochs, augmentation) against 2.6\% ``reported for the baseline in its original paper'' (step schedule, 90 epochs). Write the decomposition \eqref{eq:controls:decomp} for this comparison, say which terms the paper has measured, and design the $2\times2$ experiment that measures the others. (The description is invented.)\label{ex:controls:critique2}

\item \research (When can a search beat random by more than noise?) \Cref{thm:controls:gauss} describes random search; adaptive searches can exploit locality, the similarity of neighbors in the space, which \textcite{YingEtAl2019} measured in NAS-Bench-101. Open: a model of a space with locality, a proxy with noise and a fixed bias, and an adaptive strategy, in which the advantage over random search with the same proxy can be computed, with conditions under which it exceeds the noise of one final training.\label{ex:controls:research}
\end{exercises}

\notesandrefs
Architecture search grew out of hyperparameter optimization and neuro-evolution; the survey of \textcite{ElskenMetzenHutter2019} organizes it, as this chapter does, by search space, search strategy and performance estimation, our proxy, and \textcite{ZelaEtAl2020} explained and repaired the failures of DARTS.

\Textcite{LindauerHutter2020} turned the lessons of the reality checks into a checklist, and the tabular benchmarks NAS-Bench-101 \parencite{YingEtAl2019} and NAS-Bench-201 \parencite{DongYang2020} made the controls nearly free. Measuring the space itself is the idea of \textcite{RadosavovicEtAl2019}, who compare design spaces by the error distributions of models sampled from them: the placebo arm promoted from a baseline to an instrument.

The vocabulary of controls is borrowed from clinical trials \parencite{ICH2000,SacksChalmersSmith1982} and epidemiology \parencite{LipsitchEtAl2010}, and selection on a noisy measurement is old in quantitative genetics \parencite{FalconerMackay1996}.

\paragraph{Translating the literature.} What \textcite{LindauerHutter2020} call random sampling and random search are our placebo and sham arms. The ``random policy'' of \textcite{YuEtAl2020} is our placebo repeated over seeds, and their $p(>\mathrm{random})$ is the probability that a searched architecture beats one random architecture. The relative improvement $100(\mathrm{Acc}_m-\mathrm{Acc}_r)/\mathrm{Acc}_r$ of \textcite{YangEsperancaCarlucci2020} is our gain $G$ in accuracy, in percent of the placebo's accuracy. The test regret of \textcite{YingEtAl2019} is our $\mu(A_{\cS})-\mu^\star$ in accuracy, and their epoch budget is our proxy fidelity $f$. In the breeder's equation $R=h^2S$ of \textcite{FalconerMackay1996}, $R$ is our true gain, $S$ the selection differential of the proxy ($\sigma_Pe_k$ for the best of $k$) and $h^2$ our $\rho^2$ in the additive model; their selection intensity $i$ is $e_k$ when one of $k$ is kept.

\printchapterbib
\end{refsection}
